\PassOptionsToPackage{sols}{handout}
\documentclass[main]{subfiles}
\setcounterrefbeforedot{chapter}{m-ws6-3}
\setcounterrefafterdot{section}{m-ws6-3}
\addtocounter{section}{-1}
\usepackage{solutions}

\begin{document}

\ifSubfilesClassLoaded{%
  \chaptermark{\chaptersix}
}{}

\section{Integration by Parts} \label{ws6-3}

\begin{problemonly}
    \begin{framed}

\textbf{Key Ideas}:

\begin{itemize}[topsep=0pt,partopsep=0pt,parsep=8pt,itemsep=0pt]


  \item You should be comfortable with integration by parts as a
  way of ``undoing'' the Product Rule.

\item You should begin to develop flexibility in figuring out ``what to
  pick for $u$ and $dv$.''  There is a lot of trial and error involved in this,
  so don't feel discouraged if your first guesses don't work out: that's
  normal!  With practice, you will develop stronger intuition about what
  to pick.

\item You should be comfortable using integration by parts on
  both definite and indefinite integrals, and you should use appropriate notation when dealing with the endpoints of definite integrals

  \item You should keep in mind that you can always check whether you've correctly found an antiderivative $F$ of a function $f$ by differentiating $F$.

  \item You should be able to use integration by parts to evaluate $\displaystyle \int \ln(x)\;dx$. 


\end{itemize}
\end{framed}
\end{problemonly}

\begin{enumerate}
\item Recap! 

\begin{problem}[\vfill]
    What were the two main uses of substitution when evaluating integrals?
\end{problem}

\begin{solution}
    To undo the chain rule and to make the integrand simpler.
\end{solution}

\item Let's try to evaluate the integral $\displaystyle \int x \cdot \sin(x) \;dx$.
\begin{enumerate}
    \item 
    \begin{problem}[\vfill]
    Is it a basic antiderivative? 
\end{problem}

    \begin{solution}
        No, this is not one of our basic antiderivatives.
    \end{solution}

    \item 
    \begin{problem}[\vfill]
        Can we use algebra to simplify or change the integrand? 
    \end{problem}

    \begin{solution}
        No, there isn't any algebraic move available to reduce or simplify the integrand.
    \end{solution}

    \item 
    \begin{problem}[\vfill]
        Can we use substitution? Why or why not?
    \end{problem}

    \begin{solution}
        There are only a limited number of possibilities for $u$ and none of them work.
    \end{solution}
\end{enumerate}

Integration by substitution allowed us to ``undo" the chain rule. In this lesson, we are going to learn how to ``undo" the product rule!

\pspace{\newpage}

\item
\begin{enumerate}
    \item
    \begin{problem}[\vfill]
        State the Product Rule for the derivative of
        the product of the functions $u(x)$ and $v(x)$.
    \end{problem}

    \begin{solution}
        \begin{align*}
        \frac{d}{dx}\!\left[u(x)v(x)\right] &= u'(x)v(x) + u(x)v'(x) 
        % \\
        % \intertext{Or, using Leibniz notation:}
        % & = v(x)\frac{d}{dx}\!\left[u(x)\right] + 
        % u(x)\frac{d}{dx}\!\left[v(x)\right]
        \end{align*}
    \end{solution}

    \item
    \begin{problem}[\vfill]
        Integrate both sides with respect to $x$.
    \end{problem}

    \begin{solution}
        \begin{align*}
        \int \frac{d}{dx}\!\left[u(x)v(x)\right]  dx &= \int \left[u'(x)v(x) + u(x)v'(x) \right] dx
        \end{align*}
        By the Fundamental Theorem of Calculus, the left-hands side evaluates to $u(x)v(x)$, while we can use the properties of integrals to rewrite the right-hand side and have
        \begin{align*}
                    u(x)v(x)  &= \int \left[u'(x)v(x)\right]dx + \int \left[u(x)v'(x) \right]dx
        \end{align*}
    \end{solution}

    \item 
    \begin{problem}[\vfill]
        Rewrite your answer to part (b) in Leibniz notation.
    \end{problem}

    \begin{solution}
        Our solution to Part (b) can be written in Leibniz notation as 
        \begin{align*}
            uv & = \int v \frac{du}{dx}\; dx + \int u \frac{dv}{dx} \;dx
        \end{align*}
    \end{solution}
    
\end{enumerate}


\item We are going to use this result to calculate $\displaystyle \int x \cdot \sin(x) \;dx$.
\begin{enumerate}

\item\label{formula}\begin{problem}[\vfill] Rearrange the formula you found in \#3c to find $\displaystyle \int u \frac{dv}{dx}  \, dx$ in terms of  $\displaystyle \int v  \frac{du}{dx}  \, dx$. \end{problem} 

\begin{solution} 
  $$ \displaystyle \int u\frac{dv}{dx}dx =uv  - \displaystyle\int v \frac{du}{dx} dx$$
  \end{solution} 

\item \begin{problem}[\vfill] The above formula allows you to ``undo'' the product rule. If we want to calculate  $\displaystyle \int x \cdot \sin x  \, dx$, what should we make $u$ and what should make $\dfrac{dv}{dx}$? Why? 
\end{problem} 

\begin{solution} 
If we let $u=x$, then $\dfrac{du}{dx} = 1$. This will make things easier for us. So we should choose $\dfrac{dv}{dx} = \sin(x)$. This means that $v = - \cos(x)$.
\end{solution}  
\item \begin{problem}[\vfill]Use the formula from \ref{formula} and your choices for $u$ and $v$ to calculate $\displaystyle \int x \cdot \sin x   \, dx$.\end{problem} 

\begin{solution} 
  $$ \displaystyle \int u\frac{dv}{dx}dx =uv  - \displaystyle\int v \frac{du}{dx} dx$$
  $$ \displaystyle \int x \sin(x) dx =-x \cos(x)  - \displaystyle\int \cos(x) dx$$
    $$ \displaystyle \int x \sin(x) dx =-x \cos(x)  + \sin(x)$$

\end{solution} 


\end{enumerate} 

\pspace{\newpage}

\begin{framed}
    Formula for integration by parts: $\displaystyle \int u \;dv = uv - \displaystyle \int v \;du$

    This could also be written as: $\displaystyle \int f(x) g'(x) \;dx = f(x)g(x) - \displaystyle\int g(x)f'(x) \;dx$

\end{framed}


When integrating by parts, we want to choose part of our integrand as $u$ and the other part as $dv$. There is an element of trial-and-error, so in this next problem we will develop some intuition behind how to make strategic choices for $u$ and $dv$. If your first choice doesn't work, no worries! Just try again with new choices!

\item 

\note{
    I'm going to encourage trial-and-error here. Students
    are going to ask ``how do I know what to choose?" Let them
    know this practice is to develop the intuition.}

Evaluate the following integrals.
   \begin{enumerate}
\item \begin{problem}[\vfill]
 $\displaystyle \int x \cdot \ln x   \, dx$
\end{problem}

 \begin{solution}
This is not an integral that we simply know the antiderivative of. Also, our other technique of substitution fails here since no choice of $u$ helps with simplifying this integral! This is where Integration by Parts (IBP) comes in useful! 
$$\int f(x)g'(x) \, dx = f(x)g(x) - \int g(x) f'(x) \, dx$$
How can we see $\displaystyle \int \ln x \cdot x  \, dx$ as an integral of the type $\displaystyle \int f(x)g'(x) \, dx$? Let's try matching functions in the order in which they are showing up here! 

$$f(x) = \ln x \quad \quad g'(x) \, dx = x \, dx$$
Then $$f'(x) = \frac{1}{x} \quad \quad g(x) = \frac{x^2}{2}$$

So we have that:
\begin{align*}
\int \ln x \cdot x  \, dx &= \ln x \cdot \frac{x^2}{2} - \int  \frac{x^2}{2} \cdot \frac{1}{x} \, dx \\
&=  \ln x \cdot \frac{x^2}{2} - \int \frac{x}{2} \, dx \\
&= \ln x \cdot \frac{x^2}{2} - \frac{x^2}{4} + C
\end{align*}

Let's check our answer! 
\begin{align*}
\frac{d}{dx} \left(\ln x \cdot \frac{x^2}{2} - \frac{x^2}{4} + C\right) &= \frac{1}{x}\cdot  \frac{x^2}{2} + \ln x \cdot x - \frac{x}{2} + 0 \\
&= \frac{x}{2} +  \ln x \cdot x - \frac{x}{2} \\
&=  \ln x \cdot x
\end{align*}
That works! 
\end{solution}

  \item \label{ibp-basic-1}
  \begin{problem} [\vfill]
      $\displaystyle \int xe^x\,dx$
    \end{problem}

    \begin{solution} 
      Let \sub{$u = x$}, since it simplifies on differentiation, and \subd{$dv = e^x \,dx$} since we know the anti-derivative, and it isn't much more difficult than $dv$. Then \sub{$du = 1 dx$} and   \subd{$v = e^x $}:  
     So,
      \begin{align*}
        \int \sub{x} \subd{e^x} \, \subd{dx}  & =  \sub{x}\subd{e^x} - \int \subd{e^x}\, \sub{dx} \\
        & =   \boxed{xe^x - e^x + C}
      \end{align*}      
    \end{solution}
    
     \item 
     \note{This will be their first attempt at IBP with a definite integral.}
     \label{substitution-practice-1-3}

  \begin{problem} [\vfill\newpage]
    $\displaystyle \int_0^1 xe^x\,dx$
  \end{problem}

  \begin{solution}
    To compute definite integrals that would involve IBP, we start by first computing the indefinite integral, and then plugging in the bounds. In this case, from part \ref{ibp-basic-1}, we already know that $$\int xe^x\,dx = xe^x - e^x + C$$
    So:
    \begin{align*}
     \int_0^1 xe^x\,dx&= \left. (xe^x - e^x) \right|_0^1 \\
     &= (e-e) - (0-e^0)\\
     &= \boxed{1}
    \end{align*}
    \end{solution}
\end{enumerate}

\item \textit{Summary!}

    \begin{enumerate}
        \item 
        \begin{problem}
            During your work on Problem 5, what were some successful choices of $u$? What about some unsuccessful choices?
        \end{problem}

        \begin{solution}
            Student responses will vary, but this conversation is steering them towards the intuition that we want to choose $u$ so that it becomes simpler when we differentiate.
        \end{solution}

        \item 
        \begin{problem}
            Looking back to the formula for integration by parts and our work so far, what happens to our choice of $u$? How could this guide our initial choice for $u$?
        \end{problem}

        \begin{solution}
            We want to choose a part of the integrand as $u$ so that it becomes simpler through differentiating.
        \end{solution}

           \item 
        \begin{problem}
            What about our choices of $\dfrac{dv}{dx}$ (or $g'(x)dx$ if you used Lagrange notation)? What did we notice happens to our choices in the formula? How could this guide our choice?
        \end{problem}

        \begin{solution}
         We want to choose $\dfrac{dv}{dx}$ so that it does not become more complicated when we integrate.
        \end{solution}
        
    \end{enumerate}

\item 
Evaluate the following integrals.

\begin{enumerate}
    
\item
  \begin{problem} [\vfill]
    $\displaystyle \int  x^2e^x\,dx$
  \end{problem}
  
  \begin{solution} 
Let $u = x^2$ and  $\frac{dv}{dx} = e^x$. Then $\frac{du}{dx} = 2x$ and $v = e^x$. Then our integration by parts formula gives 
     $$\displaystyle \int  x^2e^x\,dx = x^2 e^x - \int 2x e^x dx$$
Notice we need to do integration by parts again to calculate  $\int 2x e^x dx$. Fortunately we have already done this in the previous problem. 
     $$\displaystyle \int  x^2e^x\,dx = x^2 e^x -2(xe^x - e^x)+c$$

  \end{solution}

\item 
  \begin{problem} [\vfill]
    $\displaystyle \int \ln x \,dx$
  \end{problem}

  \begin{solution}
    This doesn't look like an IBP problem. However, we can't use substitution here, and so let's give IBP a shot anyway! Let's pick So,  let \sub{$u = \ln x$}, since it simplifies on differentiation, and \subd{$dv = dx$}. Then \sub{$du = \frac{1}{x}dx$} and   \subd{$v = x$}:  
     So,
      \begin{align*}
 \int \sub{\ln x} \,\subd{dx} &=
      \sub{\ln x} \cdot \subd{x} - \int \subd{x} \sub{\frac{1}{x}} \, \sub{dx}\\
      &= x \ln x - \int dx \\
            & = \boxed{x \ln x - x + C}
    \end{align*}
  
  \end{solution}


\item 
  \begin{problem} [\vfill] 
    $\displaystyle \int_{0}^{\ln 3} x^2 e^{2x}\,dx$
  \end{problem}

  \begin{solution}
  Let's focus on computing $\displaystyle \int x^2 e^{2x}\,dx$ first. 
     Let \sub{$u = x^2$}, since it simplifies on differentiation, and \subd{$dv = e^{2x} \,dx$} since we know the anti-derivative, and it isn't much more difficult than $dv$. Then \sub{$du = 2x dx$} and   \subd{$v = \frac{e^{2x}}{2} $}:  
     So,
      \begin{align*}
        \int \sub{x^2} \subd{e^{2x}} \, \subd{dx}  & =  \sub{x^2}\cdot \subd{\frac{e^{2x}}{2}} - \int \subd{\frac{e^{2x}}{2}} \sub{2x}\, \sub{dx} \\
        & =   \frac{x^2e^{2x}}{2} - \int xe^{2x} \, dx
      \end{align*}    
      But to compute $\displaystyle  \int xe^{2x} \, dx$ we need IBP again! Deep breaths. We got this. Let \sub{$u = x$}, since it simplifies on differentiation, and \subd{$dv = e^{2x} \,dx$} since we know the anti-derivative, and it isn't much more difficult than $dv$. Then \sub{$du = 1dx$} and   \subd{$v = \frac{e^{2x}}{2} $}:  
       \begin{align*}
        \int \sub{x} \subd{e^{2x}} \, \subd{dx}  & =  \sub{x}\cdot \subd{\frac{e^{2x}}{2}} - \int \subd{\frac{e^{2x}}{2}}\, \sub{dx} \\
        & =   \frac{xe^{2x}}{2} - \frac{1}{2}\int e^{2x} \, dx \\
        &=  \frac{xe^{2x}}{2} - \frac{e^{2x}}{4} + C
      \end{align*}    
      
     Finally, we have that $$ \int x^2 e^{2x}\,dx = \frac{x^2e^{2x}}{2} - \left(\frac{xe^{2x}}{2} - \frac{e^{2x}}{4} + C\right) = \frac{x^2e^{2x}}{2} -\frac{xe^{2x}}{2} + \frac{e^{2x}}{4} -C$$
     So:
     \begin{align*}
     \int_{0}^{\ln 3} x^2 e^{2x}\,dx &= \left. \left( \frac{x^2e^{2x}}{2} -\frac{xe^{2x}}{2} + \frac{e^{2x}}{4} \right) \right|_0^{\ln3} \\
     &= \left( \frac{(\ln 3)^2e^{2\ln 3}}{2} -\frac{\ln 3e^{2\ln 3}}{2} + \frac{e^{2\ln 3}}{4} \right) -  \frac{e^{0}}{4}  \\
     &= \frac{(\ln 3)^2e^{2\ln 3}}{2} -\frac{\ln 3e^{2\ln 3}}{2} + \frac{e^{2\ln 3}}{4} - \frac{1}{4} \\
     &= \boxed{\frac{9}{2} (\ln3)^2 - \frac{9}{2}\ln 3 + 2}
     \end{align*}
  \end{solution}  
  
\end{enumerate}


\end{enumerate}



\end{document}